\par\noindent\textbf{(1)} 覆盖空气和介质的两块区域承受相同电压，因此并联电容 \(C(x)=\epsilon_0w(L-x)/d+\epsilon_0\epsilon_rw x/d=\epsilon_0w[L+(\epsilon_r-1)x]/d\)。其导数 \(C'(x)=\epsilon_0w(\epsilon_r-1)/d>0\)。
\par\noindent\textbf{(2)} 若电容器与电源断开、总电荷 \(Q\) 固定，储能 \(U_Q=Q^2/(2C(x))\)。机械力为 \(F_Q=-dU_Q/dx=Q^2C'(x)/(2C(x)^2)>0\)，介质被吸入；力随 \(x\) 变化，因为电压随电容变化。
\par\noindent\textbf{(3)} 若端电压被保持 \(V\)，电容场能 \(U_C=C(x)V^2/2\) 随插入而增大，不能直接取其负导数求力。包括电源功后的热力学势 \(\mathcal F=U_C-VQ=-C(x)V^2/2\)，故 \(F_V=-d\mathcal F/dx=V^2C'(x)/2>0\)。
\par\noindent\textbf{(4)} 在同一瞬时状态 \(Q=CV\)，代入 \(F_Q=Q^2C'/(2C^2)\) 得 \(F_Q=V^2C'/2=F_V\)。两种协议的力函数不同，是因为移动后 \(Q\) 或 \(V\) 的变化规则不同；忽略电源会导致错误的力符号。
